\honest{The Halo Effect}{Eliezer Yudkowsky}{2007}

The affect heuristic, the subject of my post three days ago, is a general feeling of good or bad spilling into other judgments. The halo effect is the same thing in judgments of people. For the evidence I quote Robert Cialdini's \textsc{Influence} at length.

Cialdini says that we credit good-looking people with ``talent, kindness, honesty, and intelligence,'' citing a review by Eagly and colleagues. \nb{That review found the effect near zero for integrity and concern for others, and moderate on average.} Attractive candidates in a Canadian election got more than two and a half times as many votes, and voters denied being influenced. \nb{Cialdini's notes give the source as an unpublished manuscript. Later published work finds the same direction on a smaller scale.} Grooming beat qualifications in a simulated interview, attractive workers earn 12 to 14 per cent more, and in a Pennsylvania study attractive defendants were ``twice as likely to avoid jail.'' \nb{That is a raw comparison. The attractive defendants also faced less serious charges, and the study's author called the relationship weak.} Attractive people also get more help and persuade more.

The effect of looks on judgments of honesty, I say, is ``a clear example of bias,'' since the two have no legitimate connection. But honesty and intelligence may really go together: ``Finding the truth, and saying the truth, are not as widely separated in nature as looking pretty and looking smart.''

So, I suggest, there ``may be a similar halo effect for kindness, or intelligence.'' Be suspicious of someone who seems intelligent, honest, altruistic, kindly and serene all at once, and of a world that sorts into devils and angels. \nb{This was the original finding. Thorndike named the halo effect in 1920 for ratings of intelligence, industry, reliability and character that correlated ``too high and too even,'' and he too allowed that such traits are really correlated to some degree.}

Finally: be a little more skeptical of attractive political candidates.
